% Unit 116 terjemahan Bahasa Indonesia dari chapter8.tex baris 1566--1703.
% Sumber: Methods of Algebra, Volume 2, commit
% 9a5803ff77dd3257484cb177f851a73770a59dd3 (CC BY 4.0).
% stable-unit-id: o014.aljabr2.chapter8.bisimplicial-b
% Cakupan: Bagian 8.8 dari peta Eilenberg--Zilber hingga struktur monoidal.

% segment-id: o014.li2.chapter8.bisimplicial-b.p001
Definisi~\ref{def:bisimplicial-diagonal} memberikan funktor aditif
$d:\cate{s}^2\mathcal{A}\to\cate{s}\mathcal{A}$. Tujuan bagian ini ialah
menjelaskan, dengan asumsi $\mathcal{A}$ kategori abelian, hubungan di antara
funktor-funktor
% segment-id: o014.li2.chapter8.bisimplicial-b.q001
\[ \tot \circ \mathrm{C}^2 , \; \mathrm{C} \circ d: \; \cate{s}^2 \mathcal{A} \rightrightarrows \cate{Ch}_{\geq 0}(\mathcal{A}). \]
Hubungan ini melibatkan shuffle-$(p,q)$ $\sigma$ dari
Definisi~\ref{def:pq-shuffle} dan tandanya $\sgn(\sigma)$ dari
Definisi~\ref{def:pq-shuffle-sgn}.

% segment-id: o014.li2.chapter8.bisimplicial-b.d001
\begin{definition}\label{def:Eilenberg-Zilber}
	\index{peta Eilenberg-Zilber@peta Eilenberg--Zilber}
	\index{peta Alexander-Whitney@peta Alexander--Whitney}
	Misalkan $\mathcal{A}$ kategori aditif dan $X$ objek bisimplisial dalam
	$\mathcal{A}$. Untuk $(p,q)\in\Z_{\geq0}^2$ yang ditetapkan, tuliskan
	$n:=p+q$.
	\begin{enumerate}[(i)]
		\item Definisikan $\overline{\mathrm{EZ}}_{p,q}:X_{p,q}\to X_{n,n}$
		sebagai
		% segment-id: o014.li2.chapter8.bisimplicial-b.q002
		\[ \overline{\mathrm{EZ}}_{p, q} := \sum_{\sigma: (p,q)\text{-shuffle}}
			\sgn(\sigma) \left( \sigma_- \times \sigma_+ : [n] \times [n] \to [p] \times [q] \right)^*. \]
		Suatu argumen panjang tetapi langsung, yang tidak kita sertakan di sini,
		menunjukkan bahwa peta-peta tersebut berpadu membentuk
		morfisme kompleks rantai yang disebut \emph{morfisme shuffle} atau
		\emph{morfisme Eilenberg--Zilber}:
		% segment-id: o014.li2.chapter8.bisimplicial-b.q003
		\[ \overline{\mathrm{EZ}}: \tot(\mathrm{C}^2 X) \to \mathrm{C}(dX). \]
		\item Definisikan
		$\overline{\mathrm{AW}}_n:=\sum_{p+q=n}\overline{\mathrm{AW}}_{p,q}:
		\mathrm{C}(dX)_n\to\tot(\mathrm{C}^2X)_n$, dengan
		$\overline{\mathrm{AW}}_{p,q}:X_{n,n}\to X_{p,q}$ didefinisikan
		sebagai $(\lambda^n_p\times\rho^n_q)^*$, tempat
		% segment-id: o014.li2.chapter8.bisimplicial-b.q004
		\[\begin{array}{rl}
			\lambda^n_p: [p] \hookrightarrow [n], & i \mapsto i, \\
			\rho^n_q: [q] \hookrightarrow [n], & i \mapsto n - q + i.
		\end{array}\]
		Dengan cara serupa, peta-peta ini berpadu membentuk morfisme
		kompleks rantai
		$\overline{\mathrm{AW}}:\mathrm{C}(dX)\to\tot(\mathrm{C}^2X)$,
		yang disebut \emph{morfisme Alexander--Whitney}.
	\end{enumerate}
\end{definition}

% segment-id: o014.li2.chapter8.bisimplicial-b.p002
Morfisme-morfisme ini funktorial bukan hanya terhadap $X$, melainkan juga
terhadap kategori $\mathcal{A}$ dan funktor-funktor di antara kategori
semacam itu. Morfisme-morfisme tersebut mula-mula diperkenalkan dalam
topologi, dalam kaitannya dengan homologi hasil kali ruang.

% segment-id: o014.li2.chapter8.bisimplicial-b.p003
Selanjutnya kita perkenalkan versi ternormalisasi dari
$\overline{\mathrm{EZ}}$ dan $\overline{\mathrm{AW}}$.

% segment-id: o014.li2.chapter8.bisimplicial-b.p004
Misalkan $\mathcal{A}$ kategori abelian bermonoidal sedemikian rupa sehingga
$\otimes:\mathcal{A}\times\mathcal{A}\to\mathcal{A}$ aditif dalam setiap
variabel. Ambil $A,B\in\Obj(\cate{s}\mathcal{A})$. Karena
$\mathrm{C}^2(A\boxtimes B)=\mathrm{C}A\boxtimes\mathrm{C}B$, kita punya
% segment-id: o014.li2.chapter8.bisimplicial-b.q005
\[ \mathrm{C}A \otimes \mathrm{C}B := \tot\left( \mathrm{C}A \boxtimes \mathrm{C}B \right) = \tot \, \mathrm{C}^2 (A \boxtimes B). \]
Kita juga menggunakan notasi yang disederhanakan
% segment-id: o014.li2.chapter8.bisimplicial-b.q006
\begin{gather*}
	A \otimes B := d(A \boxtimes B), \\
	\mathrm{N}A \otimes \mathrm{N}B := \tot\left( \mathrm{N}A \boxtimes \mathrm{N}B \right).
\end{gather*}

% segment-id: o014.li2.chapter8.bisimplicial-b.l001
\begin{lemma}\label{prop:EZ-aux}
	\index[sym1]{EZ@$\overline{\mathrm{EZ}}, \mathrm{EZ}$}
	\index[sym1]{AW@$\overline{\mathrm{AW}}, \mathrm{AW}$}
	Tinjau kategori abelian bermonoidal $\mathcal{A}$ di atas. Ingat bahwa
	untuk setiap $Y\in\Obj(\cate{s}(\mathcal{A}))$,
	Proposisi~\ref{prop:Dold-Kan-v} mengidentifikasikan $\mathrm{N}Y$
	dengan $\mathrm{C}Y/\text{bagian degenerat}$. Di bawah identifikasi ini,
	morfisme $\overline{\mathrm{AW}}$ dan $\overline{\mathrm{EZ}}$ masing-masing
	berfaktor secara unik melalui diagram komutatif berikut:
	% segment-id: o014.li2.chapter8.bisimplicial-b.q007
	\[\begin{tikzcd}
		\mathrm{C}A \otimes \mathrm{C}B \arrow[r, "{\overline{\mathrm{EZ}}}"] \arrow[twoheadrightarrow, d] & \mathrm{C}(A \otimes B) \arrow[r, "{\overline{\mathrm{AW}}}"] \arrow[twoheadrightarrow, d] & \mathrm{C}A \otimes \mathrm{C}B \arrow[twoheadrightarrow, d] \\
		\mathrm{N}A \otimes \mathrm{N}B \arrow[r, "\mathrm{EZ}"'] & \mathrm{N}(A \otimes B) \arrow[r, "{\mathrm{AW}}"'] & \mathrm{N}A \otimes \mathrm{N}B .
	\end{tikzcd}\]
\end{lemma}
\begin{proof}
	Untuk $\mathrm{AW}$, cukup diverifikasi, untuk semua
	$p,q\geq0$ dan $n:=p+q$, persamaan morfisme dalam $\simpDelta$
	% segment-id: o014.li2.chapter8.bisimplicial-b.q008
	\begin{align*}
		\lambda^n_p \circ \mathrm{s}^j & = \mathrm{s}^j \circ \lambda^{n+1}_{p+1}, \\
		\rho^n_q \circ \mathrm{s}^j & = \mathrm{s}^{j+n-q} \circ \rho^{n+1}_{q+1},
	\end{align*}
	dengan syarat pada $j$ masing-masing $0\leq j\leq p$ dan
	$0\leq j\leq q$. Verifikasi ini langsung.

	Untuk $\mathrm{EZ}$, misalkan $\sigma$ sebuah shuffle-$(p,q)$ dan
	$0\leq j\leq p-1$. Tinjau pembatasan
	$\overline{\mathrm{EZ}}_{p,q}$ pada $\Image(s_j)\otimes B_q$.
	Pembahasan sebelumnya mengenai shuffle-$(p,q)$ menunjukkan bahwa
	terdapat tepat satu $0\leq k<n:=p+q$ sedemikian rupa sehingga
	$\sigma_-(k)=j$ dan $\sigma_-(k+1)=j+1$; pada saat yang sama
	$\sigma_+(k+1)=\sigma_+(k)$. Ini memberikan persamaan dalam
	$\simpDelta$
	% segment-id: o014.li2.chapter8.bisimplicial-b.q009
	\begin{align*}
		\mathrm{s}^j \sigma_- \underbracket{\mathrm{d}^k \mathrm{s}^k}_{[n] \leftarrow [n]} & = \mathrm{s}^j \sigma_- : [n] \to [p-1] , \\
		\sigma_+ \mathrm{d}^k \mathrm{s}^k & = \sigma_+ : [n] \to [q] .
	\end{align*}

	Dengan menarik balik sepanjang persamaan-persamaan ini, kita melihat bahwa
	$\overline{\mathrm{EZ}}_{p,q}|_{\Image(s_j)\otimes B_q}$ terkandung
	dalam bagian yang degenerat melalui $s_k$. Kasus $0\leq j\leq q-1$
	dan $\overline{\mathrm{EZ}}_{p,q}|_{A_p\otimes\Image(s_j)}$ ditangani
	dengan cara yang sama.
\end{proof}

% segment-id: o014.li2.chapter8.bisimplicial-b.p005
Hasil berikut juga disebut teorema Eilenberg--Zilber yang diperumum.

% segment-id: o014.li2.chapter8.bisimplicial-b.th001
\begin{theorem}[S.\ Eilenberg, J.\ A.\ Zilber, P.\ Cartier]\label{prop:Eilenberg-Zilber}
	\index{teorema Eilenberg-Zilber@teorema Eilenberg--Zilber}
	Misalkan $\mathcal{A}$ kategori abelian dan
	$X\in\Obj(\cate{s}^2\mathcal{A})$. Morfisme kanonik yang didefinisikan
	di atas,
	% segment-id: o014.li2.chapter8.bisimplicial-b.q010
	$\begin{tikzcd}
		\tot \left(\mathrm{C}^2(X)\right) \arrow[shift left, r, "\overline{\mathrm{EZ}}"] & \mathrm{C}(d(X)) \arrow[shift left, l, "\overline{\mathrm{AW}}"]
	\end{tikzcd}$
	saling invers pada tingkat homologi; khususnya, keduanya
	kuasi-isomorfisme.
\end{theorem}
\begin{proof}
	% Reference: Kerodon, Proof of Theorem 2.5.7.14
	Karena $\cate{s}^2\mathcal{A}\simeq\cate{s}(\cate{s}\mathcal{A})$ dan
	$\cate{Ch}^2_{\geq0}(\mathcal{A})\simeq
	\cate{Ch}_{\geq0}(\cate{Ch}_{\geq0}(\mathcal{A}))$, korespondensi
	Dold--Kan memiliki versi ganda, yang tetap kita tulis sebagai ekuivalensi
	kategori
	$\Gamma:\cate{Ch}^2_{\geq0}(\mathcal{A})\to\cate{s}^2\mathcal{A}$.
	Karena itu, kita boleh mengandaikan $X=\Gamma A$ untuk
	$A\in\Obj(\cate{Ch}^2_{\geq0}(\mathcal{A}))$. Jika
	$p+q>n$ mengakibatkan $A_{p,q}=0$, kita katakan bahwa $A$ berderajat
	$\leq n$.

	Mula-mula andaikan $\mathcal{A}=\cate{Ab}$, dengan struktur monoidal
	dari hasil kali tensor $\otimes=\otimes_{\Z}$. Perhatikan bahwa
	$\tot$, $\mathrm{C}^2$, $\mathrm{C}$, dan $d$ semuanya funktor aditif
	eksak kanan dan mempertahankan jumlah langsung sebarang. Morfisme
	kanonik $\overline{\mathrm{EZ}}$ dan $\overline{\mathrm{AW}}$ juga
	kompatibel dengan jumlah langsung sebarang. Tuliskan $A$ sebagai
	gabungan terfilter subobjek-subobjek hingga. Karena kolimit
	terfilter $\varinjlim$ mempertahankan homologi, masalah tereduksi pada
	kasus $A$ berderajat $\leq n$ untuk suatu $n\geq0$ tetap. Kita
	berargumen dengan induksi pada $n$.

	Trunkasi brutal pada $A$ memberikan barisan eksak pendek
	$0\to A'\to A\to A''\to0$, dengan
	\begin{compactitem}
		\item $A'$ berderajat $\leq n-1$;
		\item semua suku tak nol $A''$ terletak pada bagian $p+q=n$;
		\item barisan eksak pendek itu terbelah pada setiap biderajat $(p,q)$.
	\end{compactitem}
	Deskripsi funktor $\Gamma$ memperlihatkan bahwa
	$0\to\Gamma A'\to\Gamma A\to\Gamma A''\to0$ juga barisan eksak pendek
	yang terbelah suku demi suku dalam $\cate{s}^2\cate{Ab}$. Secara ketat,
	di sini diperlukan versi ganda $\Gamma$, tetapi bentuknya serupa.

	Dengan demikian diperoleh diagram komutatif dengan baris-baris eksak
	dari kompleks rantai
	% segment-id: o014.li2.chapter8.bisimplicial-b.q011
	\[\begin{tikzcd}
		0 \arrow[r] & \tot\left( \mathrm{C}^2(\Gamma A') \right) \arrow[shift left, d, "\overline{\mathrm{EZ}}"] \arrow[r] & \tot\left( \mathrm{C}^2(\Gamma A) \right) \arrow[shift left, d] \arrow[r] & \tot\left( \mathrm{C}^2(\Gamma A'') \right) \arrow[shift left, d] \arrow[r] & 0 \\
		0 \arrow[r] & \mathrm{C}(d\Gamma A') \arrow[shift left, u, "\overline{\mathrm{AW}}"] \arrow[r] & \mathrm{C}(d\Gamma A) \arrow[shift left, u] \arrow[r] & \mathrm{C}(d\Gamma A'') \arrow[shift left, u] \arrow[r] & 0.
	\end{tikzcd}\]
	Diagram ini secara induktif mereduksi masalah ke kasus $A''$. Dengan
	kata lain, di bawah hipotesis induksi, masalah semula hanya bergantung
	pada $(A_{p,q})_{p+q=n}$.

	Karena itu, kita cukup menangani kasus ketika $A$ adalah jumlah
	langsung beberapa salinan
	$C\otimes(\mathrm{N}(\Z\Delta^p)\boxtimes
	\mathrm{N}(\Z\Delta^q))$, dengan $p+q=n$ dan
	$C\in\Obj(\cate{Ab})$. Selanjutnya kita dapat mereduksinya ke satu
	salinan untuk satu pasangan $(p,q)$, lalu mengeluarkan $\identity_C$
	dari semua pemetaan untuk mereduksi pada kasus khusus $C=\Z$.
	Perhatikan bahwa bagian $\mathrm{N}(\Z\Delta^p)$ pada derajat kurang dari
	$p$ tidak memengaruhi masalah; hal yang sama berlaku dengan $q$ sebagai
	pengganti $p$. Sejalan dengan itu,
	$X=\Gamma A\simeq\Z(\Delta^p\times\Delta^q)$. Berdasarkan
	Teorema~\ref{prop:Dold-Kan-N-C} dan Lema~\ref{prop:EZ-aux}, cukup
	dibuktikan bahwa
	% segment-id: o014.li2.chapter8.bisimplicial-b.q012
	\begin{equation}\label{eqn:Eilenberg-Zilber-aux}
		\begin{tikzcd}
			\mathrm{N}(\Z\Delta^p) \otimes \mathrm{N}(\Z\Delta^q) \arrow[shift left, r, "\mathrm{EZ}"] & \mathrm{N}(\Z(\Delta^p \times \Delta^q)) \arrow[shift left, l, "\mathrm{AW}"]
		\end{tikzcd}
		\; \text{saling invers pada tingkat homologi.}
	\end{equation}

	Poset $[p]$, $[q]$, dan $[p]\times[q]$ masing-masing memiliki elemen
	terkecil $0$ atau $(0,0)$, sehingga ekuivalensi homotopi $\mathrm{N}i$
	dan $\mathrm{N}q$ dari Contoh~\ref{eg:homology-poset-e} berlaku. Mudah
	diperiksa bahwa diagram berikut komutatif:
	% segment-id: o014.li2.chapter8.bisimplicial-b.q013
	\begin{equation*}\begin{tikzcd}
		\Z \otimes \Z \arrow[d, "{\mathrm{N}i \otimes \mathrm{N}i}"'] \arrow[r, "\sim", "\text{standar}"'] & \Z \arrow[d, "{\mathrm{N}i}"] \arrow[r, "\sim", "\text{standar}"'] & \Z \otimes \Z \arrow[d, "{\mathrm{N}i \otimes \mathrm{N}i}"'] \\
		\mathrm{N}(\Z\Delta^p) \otimes \mathrm{N}(\Z\Delta^q) \arrow[r, "\mathrm{EZ}"'] & \mathrm{N}(\Z(\Delta^p \times \Delta^q)) \arrow[r, "\mathrm{AW}"'] & \mathrm{N}(\Z\Delta^p) \otimes \mathrm{N}(\Z\Delta^q).
	\end{tikzcd}\end{equation*}
	Baris pertama terkonsentrasi pada derajat nol, dan isomorfisme standar
	$\Z\otimes\Z\simeq\Z$ memetakan $1\otimes1$ ke $1$. Verifikasi
	komutativitas hanya melibatkan suku berderajat nol dari $\mathrm{EZ}$
	dan $\mathrm{AW}$, sehingga langsung. Ini membuktikan
	\eqref{eqn:Eilenberg-Zilber-aux}.

	Sekarang kita kembali ke kategori $\mathcal{A}$ yang lebih umum. Andaikan terdapat funktor
	setia eksak (Proposisi~\ref{prop:faithful-exact})
	$F:\mathcal{A}\to\cate{Ab}$. Secara definisi, $\tot$, $\mathrm{C}^2$,
	$\mathrm{C}$, dan $d$ komutatif dengan $F$, sedangkan
	$\overline{\mathrm{EZ}}$ dan $\overline{\mathrm{AW}}$ jelas kompatibel
	dengan $F$. Sifat setia dan eksak lalu mereduksi masalah ke kasus yang
	telah diketahui, yaitu $\mathcal{A}=\cate{Ab}$.

	Jika $\mathcal{A}$ memiliki generator projektif, misalnya
	$\mathcal{A}=\cated{Mod}R$ untuk suatu gelanggang $R$, selalu ada
	funktor setia eksak menuju $\cate{Ab}$. Untuk kategori abelian umum
	$\mathcal{A}$, keberadaan $F$ dijamin oleh Teorema Freyd--Mitchell
	\ref{prop:FM}; untuk itu, semesta Grothendieck dapat diperbesar
	seperlunya agar $\mathcal{A}$ menjadi kategori kecil.
\end{proof}

% segment-id: o014.li2.chapter8.bisimplicial-b.r001
\begin{remark}[Struktur monoidal]
	Misalkan $\mathcal{A}$ kategori monoidal dan $\otimes$ aditif dalam
	setiap variabel. Lema~\ref{prop:EZ-aux} telah memberikan
	morfisme kanonik berbentuk
	% segment-id: o014.li2.chapter8.bisimplicial-b.q014
	\[\begin{tikzcd}[row sep=tiny]
		\mathrm{C}A \otimes \mathrm{C}B \arrow[r, "{\overline{\mathrm{EZ}}}"] & \mathrm{C}(A \otimes B) \arrow[r, "{\overline{\mathrm{AW}}}"] & \mathrm{C}A \otimes \mathrm{C}B, \\
		\mathrm{N}A \otimes \mathrm{N}B \arrow[r, "\mathrm{EZ}"'] & \mathrm{N}(A \otimes B) \arrow[r, "{\mathrm{AW}}"'] & \mathrm{N}A \otimes \mathrm{N}B
	\end{tikzcd}\]
	untuk $A,B\in\Obj(\cate{s}\mathcal{A})$. Kasus khusus teorema
	Eilenberg--Zilber yang diperumum untuk $X=A\boxtimes B$ menyatakan
	bahwa morfisme-morfisme tersebut saling invers pada tingkat homologi.
	Ingatlah bahwa $\cate{Ch}_{\geq0}(\mathcal{A})$ juga kategori monoidal
	(Contoh~\ref{eg:tensor-bisimplicial}), dan terdapat morfisme yang jelas
	% segment-id: o014.li2.chapter8.bisimplicial-b.q015
	\[\begin{tikzcd}[column sep=large]
		\munit \;\text{pada derajat nol} \arrow[r, "\sim"] & \mathrm{N}(\mathrm{const}(\munit)) \arrow[shift left, r, "\text{inklusi}"] & \mathrm{C}(\mathrm{const}(\munit)). \arrow[shift left, l, "\text{hasil bagi bagian degenerat}"]
	\end{tikzcd}\]

	Dari sudut pandang abstrak, morfisme shuffle (atau morfisme
	Alexander--Whitney) memberikan struktur funktor
	monoidal longgar kanan (atau longgar kiri) pada $\mathrm{C}$ dan
	$\mathrm{N}$; lihat \cite[Catatan 3.1.8]{Li1}. Untuk morfisme shuffle,
	secara kasar ini berarti bahwa $\overline{\mathrm{EZ}}_{p,q}$ dan
	$\mathrm{EZ}_{p,q}$ memenuhi asosiativitas dan kompatibel dengan unit;
	kasus $\overline{\mathrm{AW}}$ dan $\mathrm{AW}$ bersifat dual.
	Pernyataan-pernyataan ini pun dapat diverifikasi secara langsung,
	meskipun panjang. Misalnya, asosiativitas morfisme shuffle melibatkan
	\eqref{eqn:shuffle-sign-1}, sedangkan kasus morfisme Alexander--Whitney
	lebih sederhana.

	Jika $\mathcal{A}$ kategori monoidal simetris,
	$\overline{\mathrm{EZ}}$ dan $\mathrm{EZ}$ juga mempertahankan struktur
	braid Koszul yang diperkenalkan dalam \eqref{eqn:Koszul-braiding}.
	Untuk $\overline{\mathrm{EZ}}$, hal ini setara dengan komutativitas
	diagram berikut dalam $\mathcal{A}$:
	% segment-id: o014.li2.chapter8.bisimplicial-b.q016
	\[\begin{tikzcd}
		(\mathrm{C}A)_p \otimes (\mathrm{C}B)_q \arrow[r, "{\overline{\mathrm{EZ}}_{p, q}}"] \arrow[d, "{(-1)^{pq} \mathrm{swap}}"'] & \mathrm{C}(A \otimes B)_{p+q} \arrow[d, "{\mathrm{C}(\mathrm{swap})}"] \\
		(\mathrm{C}B)_q \otimes (\mathrm{C}A)_p \arrow[r, "{\overline{\mathrm{EZ}}_{q, p}}"'] & \mathrm{C}(B \otimes A)_{q+p}.
	\end{tikzcd}\]
	Di sini $p,q\in\Z_{\geq0}$. Peta $\mathrm{swap}$ pada kedua kolom
	berasal dari struktur braid
	$c(X,Y):X\otimes Y\rightiso Y\otimes X$ pada $\mathcal{A}$ beserta
	aksinya yang terinduksi suku demi suku pada objek simplisial; peta
	itu sendiri tidak melibatkan tanda. Bukti langsung mengikuti definisi
	dan \eqref{eqn:shuffle-sign-0}. Sebaliknya,
	$\overline{\mathrm{AW}}$ dan $\mathrm{AW}$ tidak selalu kompatibel
	dengan struktur braid.

	Secara garis besar, funktor korespondensi Dold--Kan $\mathrm{N}$ dari
	Teorema~\ref{prop:Dold-Kan}, beserta versi tak ternormalisasinya
	$\mathrm{C}$, bersifat monoidal baik dalam arti longgar kanan maupun
	longgar kiri. Data yang diperlukan masing-masing diberikan oleh
	morfisme shuffle dan morfisme Alexander--Whitney. Jika $\mathcal{A}$
	kategori monoidal simetris, struktur funktor monoidal longgar kanan itu
	juga simetris. Konsep funktor monoidal bilonggar
	\cite[Chapters 3, 5]{AM10} memberikan perumusan yang lebih tepat bagi
	struktur-struktur tersebut.

	% Morfisme shuffle juga memberi penjelasan a priori bagi braid Koszul:
	% swap pada objek simplisial tidak melibatkan tanda, dan melalui
	% korespondensi Dold--Kan ia menjadi braid Koszul pada kompleks rantai.
\end{remark}
